\bookchapter{الفصل 23: التفكير البرمجي ومبادئ الأتمتة}

\chapterintro{في الفصلين السابقين، بنينا “جدار الحدود” الآمن للهاتف، ورأينا “دماغ الذكاء الاصطناعي” الذي يزداد ذكاءً داخل الجدار. الآن، لنفتح أحد الإمكانات النهائية للهاتف: لنجعله “يفكر” مثل المبرمج، ويكمل تلقائيًا المهام المتتالية التي تحددها. قد تقول: “أنا لا أعرف كتابة الأكواد، فما علاقة هذا بي؟” لا تقلق، ما سنناقشه في هذا الفصل لا يتطلب منك كتابة سطر واحد من الكود — فقد خفضت الهواتف الحديثة عتبة “البرمجة” إلى مستوى السحب والإفلات.}

\booksection{أولًا: ما هي الأتمتة بالضبط؟}

(اختياري ولكنه قوي): إذا كانت الحالة A افعل X، وإلا افعل Y. على سبيل المثال “إذا كان يوم عمل فشغل بودكاست إخباري، وإذا كان عطلة نهاية الأسبوع فشغل الموسيقى”، وهذا يسمى التفريع الشرطي.\textbf{ثانيًا: “الاختصارات” في iOS — أداة الأتمتة السحرية لمستخدمي آبل}تطبيق “الاختصارات” (Shortcuts) من آبل هو تطبيق مدمج لا يحتاج إلى تنزيل. مفهومه الأساسي هو تجميع العمليات المختلفة معًا مثل لبنات البناء. لنقل بعض المفاهيم الأساسية أولًا:

“الاختصار”

: تدفق أتمتة كامل، يمكن تشغيله بالنقر اليدوي أو تفعيله صوتيًا عبر Siri. على سبيل المثال تقول “مرحبًا Siri، وضع العودة للمنزل”، فينفذ تلقائيًا سلسلة العمليات التي ضبطتها.

\begin{enumerate}
  \item \textbf{"الأتمتة"}: تعليمات تعمل تلقائيًا بشرط تشغيل، ولا تحتاج إلى تفعيل يدوي. على سبيل المثال "عندما يتصل الهاتف ببلوتوث السيارة، افتح تطبيق الخرائط تلقائيًا". تُضبط الأتمتة في علامة تبويب "الأتمتة"، وهي منفصلة عن "الاختصارات".
  \item \textbf{"الإجراء"}: أصغر وحدة تنفيذ. فتح تطبيق هو إجراء، وإرسال رسالة هو إجراء، والحصول على الطقس هو إجراء، وضبط السطوع هو إجراء. يمكنك ربط مئات الإجراءات معًا.
  \item \textbf{إليك بعض الأمثلة العملية المفيدة فعلًا، افهمها بتنفيذها مرة واحدة:}المثال الأول: اختصار "روتين الصباح" بضغطة واحدة
\end{enumerate}

\booksection{ثانيًا: “الاختصارات” في iOS — أداة الأتمتة السحرية لمستخدمي آبل}

المثال الثاني: “وضع توفير الطاقة” التلقائي

\begin{itemize}
  \item \textbf{في علامة تبويب "الأتمتة"، أنشئ أتمتة اختر "بطارية" اضبط "أقل من 30\%" أضف الإجراءات "ضبط وضع الطاقة المنخفضة على تشغيل" و"خفض سطوع الشاشة إلى 30\%" و"إيقاف تشغيل البلوتوث". من الآن فصاعدًا، عندما تنخفض بطارية الهاتف عن 30\% سيدخل تلقائيًا في وضع توفير الطاقة دون أن تقلق.}المثال الثالث: التسجيل التلقائي للحسابات بعد مسح الرمز
  \item \textbf{في كل مرة تستخدم فيها WeChat/AlPay للمسح والدفع، سيقرأ الهاتف إشعار الدفع، ويسجل تلقائيًا المبلغ والتاجر في مذكرة أو جدول Numbers. يتطلب هذا استخدام شرط التشغيل "استلام رسالة/إشعار" في "الأتمتة" ضمن الاختصارات (iOS 17+ يدعم قراءة معاملات المحفظة)، مع إجراء "الإلحاق بـ Numbers". لمن يريد تسجيل الحسابات لكنه يكسل عن التسجيل اليدوي، فهذا منقذ حقيقي.}هناك استخدام ذكي آخر لا يعرفه كثيرون: يمكنك إضافة اختصار إلى الشاشة الرئيسية ليصبح أيقونة تطبيق. على سبيل المثال، أنشئ اختصار "استدعاء سيارة للعودة للمنزل"، ضع الأيقونة على سطح المكتب، وانقر مرة واحدة ليفتح DiDi مباشرة ويضبط الوجهة على المنزل، وهو أسرع بكثير من فتح التطبيق بنفسك وإدخال العنوان ببطء.
  \item \textbf{ثالثًا: أتمتة Android — Tasker وأصدقاؤه}في جانب Android، الوضع أعقد قليلًا، لكن السقف أيضًا أعلى. لنبدأ بالبسيط:
\end{itemize}

مستوى المبتدئ: النقر التلقائي

\booksubsection{المثال الأول: اختصار “روتين الصباح” بضغطة واحدة}

المستوى المتوسط: Tasker — ملك أتمتة Android

\booksubsection{المثال الثاني: “وضع توفير الطاقة” التلقائي}

ملف التعريف (Profile)

\booksubsection{المثال الثالث: التسجيل التلقائي للحسابات بعد مسح الرمز}

المهمة (Task)

: تحدد “ماذا تفعل بعد التشغيل”. يمكن أن تتضمن المهمة الواحدة عدة إجراءات، تُنفذ بالترتيب. تشمل أنواع الإجراءات: ضبط إعدادات النظام (الصوت، السطوع، Wi-Fi، البلوتوث، وضع الطيران)، إرسال الرسائل النصية، طلبات HTTP، تنفيذ أوامر Shell، التحكم في تشغيل الوسائط، عرض مربعات حوار، وغيرها.

\booksection{ثالثًا: أتمتة Android — Tasker وأصدقاؤه}

تشغيل Wi-Fi تلقائيًا عند الوصول للمنزل

\booksubsection{مستوى المبتدئ: النقر التلقائي}

وضع الاجتماع

\booksubsection{المستوى المتوسط: Tasker — ملك أتمتة Android}

حماية العين ليلًا

\begin{itemize}
  \item \textbf{: شرط التشغيل "الوقت 22:00-06:00" المهمة "تشغيل وضع حماية العين، خفض السطوع إلى 20\%، ضبط الشاشة على الأبيض والأسود (وضع التدرج الرمادي)". الشاشة بالأبيض والأسود لها تأثير عجيب على الرغبة في تصفح مقاطع الفيديو القصيرة — بمجرد أن تصبح الشاشة رمادية، لن ترغب في التصفح.}تشغيل البودكاست تلقائيًا عند توصيل سماعات الرأس
  \item \textbf{: شرط التشغيل "توصيل سماعات الرأس السلكية/اتصال سماعات البلوتوث" المهمة "فتح تطبيق Xiaoyuzhou وبدء تشغيل أحدث حلقة". بمجرد ارتداء سماعات الرأس تستطيع الاستماع، دون فتح قفل الهاتف للبحث عن التطبيق.}إذا شعرت أن Tasker صعب الإتقان، يمكنك تجربة شقيقه المبسط:
\end{itemize}

MacroDroid

\begin{itemize}
  \item \textbf{. واجهته أكثر ودًا وقوالبه غنية، وهو مناسب للمستخدمين الذين لا يرغبون في العبث. بالإضافة إلى ذلك، تأتي العديد من أنظمة Android الصينية الآن (MIUI من شاومي وHarmonyOS من هواوي وColorOS من OPPO) بـ "مشاهد" أو "مساعد ذكي" مدمج، وإن لم يكن بقوة Tasker، إلا أنه يكفي تمامًا للأتمتة اليومية.}رابعًا: التطبيقات اليومية للتفكير الآلي
  \item \textbf{أكبر مكسب من تعلم الأتمتة ليس توفير ذلك الوقت القليل، بل}تغيير طريقة تفكيرك في المشكلات
  \item \textbf{. بمجرد أن تعتاد على نمط التفكير "إذا حدث X، فافعل Y تلقائيًا"، ستكتشف أن كمية كبيرة من العمل المتكرر في الحياة يمكن التخلص منها. على سبيل المثال:}في كل مرة تنتهي من التقاط الصور، يتم النسخ الاحتياطي تلقائيًا إلى التخزين السحابي ثم ضغط الصور المحلية إلى حجم صغير لتوفير المساحة.
  \item \textbf{كل صباح يُقرأ التقويم تلقائيًا لجدول اليوم، ويُولَّد نص ويُرسل إلى أفراد العائلة عبر WeChat.}عند استلام رسالة استلام طرد، يُضاف تلقائيًا إلى المذكرة عنصر مهام "استلام الطرد".
\end{itemize}

عند اكتمال شحن الهاتف، يُرسل تذكير صوتي لفصل الشاحن (لتجنب الإفراط في الشحن الذي يضر البطارية).\textbf{عندما يحين وقت انتهاء الدوام يوميًا، إذا كنت لا تزال بالقرب من الشركة (عبر تحديد موقع البرج)، يأتي تذكير تلقائي “حان وقت العودة للمنزل”.}هذه تبدو متقدمة، لكن في الواقع بمجرد إتقانك لمفهومي “التشغيل + الإجراء”، بالتعاون مع الاختصارات أو Tasker، يمكنك إنجازها جميعًا. المفتاح هو

\booksection{رابعًا: التطبيقات اليومية للتفكير الآلي}

— فكثير من الناس لا يعجزون عن ذلك، بل لا يفكرون فيه أصلًا.\textbf{خامسًا: دليل تجنب الأخطاء — المشكلات الشائعة في الأتمتة}الأتمتة ليست كلها سلسة، وهناك بعض الحفر جربتها قبلك:

\begin{itemize}
  \item قد لا يعمل تشغيل الموقع بدقة
  \item : يستخدم تشغيل "الوصول/المغادرة" في iOS تقنية السياج الجغرافي، وتعتمد الدقة على إشارة GPS وتحسين النظام. أحيانًا لا يعمل التشغيل إلا بعد 10 دقائق من وصولك للمنزل. يُنصح باستخدام تشغيل اتصال Wi-Fi أو البرج الخلوي بدلًا من تشغيل الموقع، فهو أكثر موثوقية وأكثر توفيرًا للطاقة.
  \item "الأتمتة" في الاختصارات تحتاج إلى إذن
  \item : في iOS، معظم عمليات التشغيل في علامة تبويب "الأتمتة" تحتاج إلى تأكيد يدوي للتشغيل (آلية أمان iOS). لكن تشغيلات اتصال Wi-Fi واتصال البلوتوث وإيقاف المنبه يمكن تشغيلها تلقائيًا بالكامل دون تأكيد.
  \item يحتاج Tasker إلى تفعيل إذن الوصول
\end{itemize}

: العديد من الوظائف المتقدمة (مثل محاكاة النقر) تحتاج إلى إذن خدمة الوصول، مما يجعل الهاتف يعرض إشعارًا دائمًا في الإعدادات. هذا تصميم أمان في Android ولا يمكن إزالته.\textbf{لا تجعل الأتمتة معقدة للغاية}: عند البدء في الأتمتة، ابدأ بتدفقات صغيرة من “2-3 إجراءات”. لا تبدأ مباشرة بتدفق فائق من 20 خطوة — فستجنّ عند تصحيحه.

\booksection{خامسًا: دليل تجنب الأخطاء — المشكلات الشائعة في الأتمتة}

اقضِ 5 دقائق في التفكير: ما هي أكثر 3 أشياء تتكرر عليك يوميًا على الهاتف؟ على سبيل المثال، الروتين الصباحي اليومي، والعادات بعد توصيل سماعات البلوتوث في كل مرة، والإعدادات قبل الخروج... ابدأ من هذه الأشياء الصغيرة عالية التكرار، واصنع أتمتة بسيطة لكل منها. حتى لو وفرت كل أتمتة 5 ثوانٍ فقط، فالمجموع مع الوقت كبير. والإشباع بعد إنجازها يدمن فعلًا.

\begin{itemize}
  \item \textbf{تدرَّب: ثلاثة تحديات أتمتة للمبتدئين}اختر الأنسب لك وأنجزه، وبعد إكماله ستكون قد دخلت الميدان:
  \item \textbf{التحدي الأول: أتمتة صباح الخير}: iOS: في تطبيق الاختصارات أنشئ اختصارًا جديدًا أضف "الحصول على الطقس الحالي" "قراءة النص" "فتح تطبيق واختيار NetEase Cloud". بعد الإكمال يمكنك قول "صباح الخير" لـ Siri للتشغيل. Android: في Tasker أنشئ مهمة أضف "الحصول على الطقس" "قراءة" "فتح تطبيق" اضبط وقت التشغيل على 7:00 يوميًا.
  \item \textbf{التحدي الثاني: توفير الطاقة تلقائيًا عند انخفاض البطارية}: iOS: علامة الأتمتة جديد "بطارية" أقل من 25\% أضف "ضبط وضع الطاقة المنخفضة على تشغيل" + "ضبط السطوع 25\%". Android: Tasker أنشئ Profile الحالة الطاقة البطارية 0-25\% المهمة "تشغيل وضع توفير الطاقة" + "خفض السطوع".
  \item \textbf{التحدي الثالث: التنقل التلقائي عند توصيل بلوتوث السيارة}: iOS: الأتمتة جديد "البلوتوث" اختر بلوتوث سيارتك أضف "فتح تطبيق واختيار خرائط AutoNavi" + "ضبط صوت الوسائط إلى 80\%". Android: Tasker Profile الحالة البلوتوث متصل اختر بلوتوث السيارة المهمة "فتح AutoNavi" + "رفع صوت الوسائط".
\end{itemize}

\textbf{ملحق الفصل 23: دليل البدء في GitHub — من استضافة الأكواد إلى منصة التعاون}\par

في الفصل السابق تحدثنا عن “التفكير البرمجي”، ولترجمة هذا التفكير إلى واقع والعمل مع آخرين في مشاريع، لا مفر من اسم — GitHub. سواء كنت طالب علوم حاسوب تشارك في مسابقات، أو تريد حفظ أعمالك البرمجية، أو تريد ببساطة رؤية ما يفعله المطورون حول العالم، فإن GitHub منصة يجب أن تعرفها. في هذا القسم سنبدأ من المفاهيم الأساسية لمساعدتك في بناء فهم واضح.

\begin{tcbpractice}{أولًا: ما هو GitHub بالضبط؟}
ببساطة، GitHub هو

\begin{enumerate}
  \item \textbf{منصة استضافة الأكواد والتعاون}قائمة على Git. أُطلق في 2008، واستحوذت عليه Microsoft في 2018، ولا يزال أكبر منصة استضافة أكواد مفتوحة المصدر في العالم، يستضيف مئات الملايين من المستودعات، ويغطي جميع لغات البرمجة الرئيسية تقريبًا.
  \item \textbf{إذا استخدمنا تشبيهًا لفهمه:}Git هو دفتر يوميات على حاسوبك، و GitHub هو منصة تدوين علنية
  \item \textbf{— يمكنك نشر أعمالك، ومتابعة مشاريع الآخرين، والمشاركة في النقاش والتعاون. إنها ليست مجرد "مستودع خلفي للمبرمجين"، بل أشبه بـ "مدينة تقنية" يبنيها المطورون حول العالم معًا — تُولد المشاريع هنا، وتتطور، وتُناقش، وتُعاد استخدامها.}بالنسبة للشخص العادي، GitHub ليس بعيد المنال أيضًا. كثير من غير المبرمجين يستخدمونه لاستضافة الملاحظات والكتب الإلكترونية والمدونات الشخصية، بل وحتى لنسخ ملفات التكوين احتياطيًا. طالما لديك أدنى حاجة إلى "تنظيم أشيائك ومشاركتها"، فإن GitHub يستحق أن تتعرف عليه.
\end{enumerate}

\end{tcbpractice}

\begin{tcbinfo}{حالة 3: أتمتة "روتين الصباح"}
\begin{itemize}
  \item \textbf{السيناريو}: كل يوم 7:00 بث تلقائي للطقس + إيقاف الكتم + فتح الأخبار.
  \item \textbf{الخطوات}: iOS: Shortcuts Automation Time trigger أضف "Get Weather" + "Set Volume" + "Open App"؛ Android: MacroDroid Trigger (Time 7:00) Actions (Media Volume/Open URL/TTS Weather).
  \item \textbf{النتيجة المتوقعة}: يمكن تحقيقه دون تثبيت، يخفض عتبة البدء.
\end{itemize}

\end{tcbinfo}
